\SourceLeaf{255}%
\Needspace{7\baselineskip}
\begin{center}
\textsc{Gauss} an \textsc{Gerling}. Göttingen, 23. Juni 1846.
\end{center}

\mbox{.\ .\ .\ .\ .\ .\ .} Die Unterscheidung der Schrauben hängt letztlich ebenso wie die Unterscheidung von Rechts und Links davon ab, dass

\begin{quote}
drei von einem Punkt $A$ ausgehende und nicht in einer Ebene liegende Geraden $AB$, $AC$, $AD$ mit drei anderen $ab$, $ac$, $ad$ in der Reihenfolge, in der sie aufgezählt sind,
\end{quote}

\begin{center}
die gleiche oder die entgegengesetzte Orientierung haben können.
\end{center}

\Needspace{14\baselineskip}
Im ersten Fall kann nämlich die Pyramide $ABCD$ allein durch allmähliche quantitative Änderungen so in die Pyramide $abcd$ übergehen, dass $A$ in $a$, $B$ in $b$ usw. übergeht. Im zweiten Fall ist ein solcher Übergang nur möglich, wenn die sich verändernde Pyramide dazwischen einmal \textls{verschwindet} ($ABCD$ in einer Ebene). Gleiche Orientierung haben

\[
\begin{array}{ccc}
AB,&AC,&AD\\
AC,&AD,&AB\\
AD,&AB,&AC,
\end{array}
\]

die entgegengesetzte Orientierung dagegen $AB$, $AD$, $AC$ usw.

Diesen \textls{Unterschied} zwischen zwei Systemen aus je drei Geraden kann man aber \textls{nicht} auf Begriffe bringen. Nur durch Bezug auf \textls{wirklich vorhandene} räumliche Dinge lässt er sich \textls{vorzeigen}. Zwei Geister können sich darüber nur verständigen, indem sie ihre Anschauungen an ein und dasselbe in der wirklichen Welt vorkommende System knüpfen. Wählt man also für das 1\textsuperscript{te} System die [Ordnung:]

\Needspace{7\baselineskip}
\begin{center}
Süd \qquad West \qquad Zenit,
\end{center}

so steht für jedes andere System aus drei Geraden fest, ob es gleich oder entgegengesetzt orientiert ist. Gleich orientiert sind etwa:

\Needspace{7\baselineskip}
\begin{center}
\begin{tabular}{cccc}
Vorne & Rechts & Oben & \\
Links & Vorne & Oben & \\
Vorne & Links & Unten & \\
Rechts & Vorne & Unten & \\
Rechts & Hinten & Oben & usw.
\end{tabular}
\end{center}
